%!TEX root = tfg.tex
\chapter{Logros principales}

\begin{abstract}
Este capítulo recoge los logros alcanzados a lo largo del proyecto, contrastándolos con los objetivos definidos al inicio del trabajo. Se evalúa el grado de cumplimiento de cada objetivo y se destacan los resultados más relevantes obtenidos.
\end{abstract}

Los objetivos planteados al inicio del proyecto han sido alcanzados en su totalidad, con la excepción del componente de validación con robot físico, que quedó fuera del alcance del trabajo por limitaciones de tiempo y disponibilidad del hardware. A continuación se detalla el grado de cumplimiento de cada objetivo.

\begin{description}
    \item[\textbf{Objetivo 1. Puesta en marcha del entorno de simulación.}] \hfill\break
    \textbf{Completado.} Se ha configurado un entorno de simulación funcional basado en Ignition Gazebo Fortress con el modelo TurtleBot4 Standard, operando sobre ROS2 Humble. El entorno reproduce fielmente el comportamiento del robot real en lo relativo a control de movimiento, lectura del LiDAR y comportamientos de seguridad del firmware del Create3. Para los experimentos dinámicos, se ha configurado un entorno Docker sobre el PC del departamento que permite ejecutar la simulación al 85--100\% de la velocidad real, lo que ha sido imprescindible para la recogida de métricas dinámicas.

    \item[\textbf{Objetivo 2. Desarrollo de nodos de control en ROS2.}] \hfill\break
    \textbf{Completado.} Se han desarrollado en C++ dos nodos de medición: un nodo estático (\texttt{data\_measure}) que captura las métricas de la trayectoria planificada sin ejecutar la navegación, y un nodo dinámico (\texttt{data\_measure\_dynamic}) que muestrea el estado del robot a 5\,Hz durante la navegación real. Ambos nodos escriben los resultados en ficheros CSV nombrados automáticamente según la configuración del algoritmo, lo que ha permitido automatizar la recogida de datos para las cuatro variantes evaluadas.

    \item[\textbf{Objetivo 3. Integración con el stack de navegación Nav2.}] \hfill\break
    \textbf{Completado.} Nav2 se ha configurado con el planificador global NavFn, el controlador local DWB y localización mediante AMCL sobre el mapa pregenerado del almacén. Se ha automatizado la inicialización de la pose mediante un \texttt{TimerAction} en el launch file, eliminando la necesidad de intervención manual al arrancar el sistema. La configuración es reproducible y está parametrizada para facilitar el cambio entre variantes de algoritmo.

    \item[\textbf{Objetivo 4. Implementación y evaluación de algoritmos de planificación.}] \hfill\break
    \textbf{Completado.} Se han evaluado cuatro configuraciones del planificador global: Dijkstra con y sin suavizado, y A* con y sin suavizado. Aunque el objetivo inicial contemplaba al menos un algoritmo alternativo adicional (RRT), su integración en Nav2 Humble habría requerido el desarrollo de un plugin personalizado, lo que quedaba fuera del alcance del trabajo. En su lugar, el análisis del suavizado de trayectorias como variable independiente ha aportado información relevante adicional sobre el comportamiento del planificador.

    \item[\textbf{Objetivo 5. Diseño del sistema de medición.}] \hfill\break
    \textbf{Completado.} El sistema de medición cubre métricas estáticas (tiempo de planificación, número de waypoints, longitud de trayectoria, irregularidad y curvatura media) y métricas dinámicas (velocidad lineal y angular, ETA calculado por Nav2, desviación respecto al plan y tiempo total de navegación). Los datos se procesan mediante un proyecto Python independiente que genera automáticamente las gráficas comparativas incluidas en el capítulo de resultados.

    \item[\textbf{Objetivo 6. Análisis comparativo.}] \hfill\break
    \textbf{Completado.} El análisis realizado permite extraer conclusiones cuantitativas y cualitativas sobre el comportamiento de cada configuración. Los resultados más relevantes son los siguientes:

\end{description}

\begin{itemize}
    \item Dijkstra produce trayectorias entre 4 y 5 veces más suaves que A* en entornos con obstáculos, medido como cambio de ángulo acumulado. Esta diferencia es estructural y no se corrige con el suavizado de NavFn.

    \item A* presenta tiempos de planificación entre un 20 y un 50\% menores que Dijkstra smooth, lo que lo hace preferible en escenarios donde la velocidad de replanificación sea crítica y la suavidad de la trayectoria sea menos relevante.

    \item El suavizado de trayectorias (\texttt{do\_refinement: true}) tiene un efecto limitado sobre la irregularidad (reducción inferior al 5\%) pero mejora moderadamente la curvatura media (entre un 10 y un 20\% en la mayoría de los goals). Su coste en tiempo de planificación es de entre 3 y 9 ms en Dijkstra.

    \item En ejecución, el controlador local DWB tiende a igualar el comportamiento de las distintas configuraciones: las diferencias de tiempo total de navegación son inferiores al 5\% y las de desviación media inferiores a 0,1 metros. La diferencia más apreciable es el perfil de velocidad angular, notablemente más errático en A* como consecuencia directa de la mayor irregularidad de sus trayectorias.

    \item Dijkstra smooth es la configuración con mejor desviación media global (1,204 m) y la que produce un comportamiento de navegación más estable y predecible, siendo la opción recomendada para entornos de almacén con obstáculos estáticos.
\end{itemize}

Más allá del cumplimiento de los objetivos, el proyecto ha permitido adquirir una comprensión profunda del ecosistema ROS2 y del stack de navegación Nav2, desde la configuración del entorno de simulación hasta el análisis cuantitativo de los resultados. El trabajo desarrollado constituye una base sólida y reproducible para futuras investigaciones en navegación autónoma de robots móviles.
